% Source: https://artofproblemsolving.com/community/c3h2579342_question_about_solving_equations
% Stable identity: aops-topic-2579342-q01
\textbf{(a)} Squaring gives $2x+3=x^2$, i.e. $(x-3)(x+1)=0$, so $x=3$ or $x=-1$. Check: $\sqrt9=3$ works, but $\sqrt{1}=1\ne-1$. So the only solution is $\boxed{x=3}$.

The step from $u=v$ to $u^2=v^2$ is not reversible: $u^2=v^2$ also allows $u=-v$. Here $x=-1$ solves $\sqrt{2x+3}=-x$ instead.

\textbf{(b)} Multiplying by $x-2$ gives $x^2=4$, so $x=\pm2$. But the original equation is undefined at $x=2$, because we multiplied by $x-2=0$ there. So the only solution is $\boxed{x=-2}$. Check: both sides equal $-1$.

In general, checking is necessary whenever a step is not reversible, such as squaring or multiplying by an expression that may be $0$. If every step is an equivalence, no solutions can be gained or lost, though checking is still a good way to catch arithmetic slips.
